\section{Divisor growth and residue currents}\label{sec:curves}

Proposition~\ref{compact:section} constructs a section with end limits
$(1,0)$ for any prescribed finite decay rate of the holomorphic
structure. Before choosing that rate, we prove an area bound with an
exponent $A$ depending only on the lifted geometry. We then choose the
holomorphic structures to decay faster than this common exponent.

Our goal is to constrain the normalizations of the resulting divisor
components. A holomorphic differential on a normalization defines a
closed current on $Y$. We place currents with exponential mass bounds
in fixed weighted cohomology spaces and prove that currents carried by
distinct curves have independent classes. The scalar degree-one group vanishes,
so scalar differentials satisfying the mass bound cannot occur. This
will force genus zero here; Section~\ref{sec:cylinder} will exclude the
remaining noncompact normalizations, using the same scalar obstruction
and the finite-dimensional cohomology for a fixed bundle $E$.

\subsection{A common area exponent}

Areas are measured with the lifted Hermitian metric; multiplicities
are included when the argument is a divisor. Put
\[
 I_l=[-l-1,-l],\qquad Y_l=h^{-1}(I_l)\quad(l\in\NN\cup\{0\}).
\]

\begin{lemma}[A geometric area exponent]\label{curve:area}
There is a number $A>0$, depending only on a fixed finite coordinate
atlas of $X$, its lifted metric, and the covering height, with the following
property.  Let $E$ be any compact datum whose holomorphic operator tends
smoothly to the trivial operator at the negative end, and let $s$ be a
nonzero holomorphic section which is bounded in the prescribed smooth
trivialization there.  If $D=(s=0)$, then
\begin{equation}\label{curve:area-bound}
 \Area(D\cap Y_l)\le C_{E,s}\exp(Al),\qquad l\ge0.
\end{equation}
The constant $C_{E,s}$ may depend on the datum, operator, and section.
The exponent $A$ is independent of their decay rate.  The same exponent
applies to every reduced component of $D$ and each of its fixed deck
translates.  Replacing $I_l$ by any fixed-width enlargement also preserves
the exponent.
\end{lemma}

\begin{proof}
Here are details of the quantitative propagation used in the proof.
Choose a compact set $Q\subset Y$ whose deck translates cover $Y$.
Choose finitely many small coordinate balls covering $Q$, and paths
joining their centers to a point $p$, together with a path from $p$ to
$T^{-1}p$.  The union of these paths is compact.  Subdividing the paths
and slightly enlarging the balls produces a finite collection of ball
chains.  Repeating its translates connects a starting ball to every
ball covering $T^{-l}Q$ in at most $N(l+1)$ transitions, for one fixed
$N$.  All balls in such a chain lie in a fixed-width enlargement of the
height interval traversed by the chain.  These statements use only
connectedness of $Y$, compactness of the chosen paths, and the identity
$h\circ T=h+1$; connectedness of individual height levels is unnecessary.

We recall the elementary analytic estimate for a transition.  If $F$ is
holomorphic on a Euclidean ball of radius $r_3$, and
$0<r_1<r_2<r_3$, the three-circle inequality is
\begin{equation}\label{curve:three-ball}
 \log\sup_{B_{r_2}}|F|
 \le \vartheta\log\sup_{B_{r_1}}|F|
       +(1-\vartheta)\log\sup_{B_{r_3}}|F|,
 \qquad
 \vartheta=\frac{\log(r_3/r_2)}{\log(r_3/r_1)}.
\end{equation}
Indeed, restricting to each complex line through the center gives the
one-variable three-circle inequality, and then taking the supremum over
the lines gives the displayed inequality.  Choose the ball chains so
that a small ball at one stage is contained in a middle ball at the next
stage, while that stage's large ball remains inside its coordinate
chart.  This is possible by subdividing the paths further.  Only
finitely many radius ratios and coordinate changes are involved.
Consequently the corresponding numbers $\vartheta$ have a positive
lower bound $\vartheta_0$.  The same construction, with a bounded number
of additional transitions, reaches the small concentric balls needed
for all the local integrations below.

On a sufficiently negative tail write the operator as $\dbar+\alpha$
in the prescribed smooth frame.  Since $\alpha\to0$ smoothly and
$\dbar\alpha=0$, on each of our large coordinate balls there is a
function $u$ with $\dbar u=\alpha$ and a uniform bound $|u|\le C_0$ on
a slightly smaller ball.  This is the local Dolbeault lemma with its
interior estimate.  One can obtain the estimate directly by the
Cauchy--Green operator in the first coordinate and then in the second:
closedness removes the first antiholomorphic derivative from the
remaining coefficient; the integrable kernels give a bound in terms
of finitely many coefficient norms on the larger polydisc.  The finite
atlas makes these bounds uniform.  Moving the start of the tail makes
the coefficient norms at most one, independently of the original
decay rate.  The frames $e^{-u}$ are holomorphic and their norms, and
the absolute values of transition functions between these frames,
are bounded above and below by geometric constants.

Let $F$ denote the coefficient of $s$ in such a holomorphic frame.
The boundedness of $s$ gives a common upper bound
$\log|F|\le U$ on the tail, where $U\ge1$ may depend on $s$.
At every transition, \eqref{curve:three-ball}, the inclusion of the
preceding small ball in the next middle ball, and the bounded frame
changes show the following.  If the logarithmic supremum on the
preceding small ball is at least $-M$, then on the next small ball it
is at least
\[
 -cM-C_1(U+1),\qquad c=\vartheta_0^{-1}>1.
\]
Increasing $c$ if needed accommodates any transition with equal radii.
Since $s$ is nonzero, unique continuation supplies a finite starting
value $M$.  Start sufficiently far down the tail that all repeated
chains stay in the region where the frame bounds hold.  Iteration over
at most $N(l+1)$ transitions gives
\begin{equation}\label{curve:log-lower}
 \sup_V\log|F|\ge -C_2 c^{Nl}
\end{equation}
on every small ball $V$ needed to cover the $l$th negative cell.
Here and below a change of the starting cell changes $C_2$, not $c$ or
$N$.  A slab $Y_l$ meets only a fixed number of these cells, because
$h$ is bounded on $Q$.

For completeness the conversion from a logarithmic supremum to an
$L^1$ bound is elementary.  Arrange nested coordinate balls of radii
$r,2r,4r$ and choose $x\in B_r$ with
$\log|F(x)|\ge-M-1$, using \eqref{curve:log-lower}.
The nonnegative function $v=U-\log|F|$ is superharmonic.  The ball
$B(x,3r)$ contains $B_{2r}$ and lies in $B_{4r}$, so its mean inequality
gives
\[
 \int_{B_{2r}}v\le \int_{B(x,3r)}v
       \le |B_{3r}|\,v(x)\le |B_{3r}|(U+M+1).
\]
The usual limiting interpretation applies at zeros.  Thus
$\int_{B_{2r}}|\log|F||\le C_3(U+M+1)$.
Choose a smooth cutoff $\chi$ equal to one on $B_r$ and supported in
$B_{2r}$.  Poincar\'e--Lelong
\cite[Chapter~III, (2.15)]{Demailly2012} and distributional integration by parts
give, up to the fixed normalization of $d^c$,
\[
 \int_D\chi\,\omega
 =c_0\int\log|F|\,dd^c(\chi\omega)
 \le C_4\int_{B_{2r}}|\log|F||.
\]
The coefficients of $dd^c(\chi\omega)$ are uniformly bounded.  No
closedness of the Hermitian form $\omega$ is needed for this local
identity.  Summing over the fixed number of balls covering a slab
and using \eqref{curve:log-lower} proves \eqref{curve:area-bound}, for
example with any fixed $A>N\log c$.  Finitely many slabs outside the
chosen tail are compact and are absorbed in $C_{E,s}$.

A reduced component has no more area than the effective divisor.
Since $T$ is an isometry and raises height by one, replacing a component
by $T^kB$ shifts the slab index by $k$.  After absorbing the finitely
many remaining slabs in the constant, its exponent is still $A$.
A fixed-width enlarged slab is contained in a bounded number of unit
slabs whose indices differ from $l$ by a bounded amount.  This proves
the final assertion as well.
\end{proof}

From now on fix this geometric $A$, and make every application of
Proposition~\ref{compact:section} with a decay rate
\begin{equation}\label{curve:gamma-choice}
 \gamma>2A+10.
\end{equation}
That proposition was proved for every prescribed finite rate, so this
is a specialization of its construction.  For each datum we obtain
$E,s,D$ with $s\to1$ at the positive end and $s\to0$ at the negative
end.  In particular $D$ is bounded above in height.  Lemma~\ref{curve:area}
applies to these choices with the already fixed exponent $A$.

\subsection{Currents and their weighted cohomology}

Let $B\subset Y$ be a closed irreducible analytic curve, possibly
noncompact and bounded above
in height, and let $f:R\to B\subset Y$ be its normalization.
Normalization is finite over each relatively compact open subset
\cite[Chapter~II, (7.12)]{Demailly2012}, so
$f$ is proper.  If $F$ is $\OO$ or one of the bundles $E$, and
$\eta$ is a holomorphic section of $K_R\otimes f^*F$, define its
pushforward current $C_\eta=f_*\eta$ by
\begin{equation}\label{curve:push-current}
 C_\eta(\varphi)=\int_R\eta\wedge f^*\varphi,
 \qquad
 \varphi\in C_c^\infty(Y,\Lambda^{0,1}\otimes F^*).
\end{equation}
The current has bidegree $(2,1)$ and coefficients in $F$.
Properness makes the integral locally finite, including at singular
points of $B$.  Stokes' theorem on $R$, applied to a compactly supported
pullback, shows $\dbar_F C_\eta=0$.
With any conformal metric on $R$ there is a fixed convention-dependent
constant $c$ such that
\begin{equation}\label{curve:mass-estimate}
 \Mass(C_\eta;V)
 \le c\int_{f^{-1}(V)}|\eta|\,|df|\,dA_R.
\end{equation}
The expression on the right is conformally invariant.  It also shows
directly that $C_\eta$ is of order zero.

We need three properties of these currents: finite-dimensional
cohomology for a fixed coefficient bundle, vanishing of scalar
cohomology in degree one, and vanishing on a positive tail of any
$(2,0)$-primitive of a $(2,1)$-current supported below a height bound. The weighted
results of Section~\ref{sec:weighted} give all three in the following
fixed spaces. The mass criterion at the end of the theorem specifies
which currents belong to them.

\begin{theorem}[Weighted cohomology for residue currents]
\label{weighted:currents}
Fix $\eps>0$ and $b>\eps$. Let $F$ be either $\OO_Y$ or an end-framed
compact datum $E$ with an integrable operator satisfying
\eqref{weighted:end-error}. Consider the current complex
\begin{equation}\label{weighted:current-complex}
 \mathcal C^{2,q}_{-\eps,b,-3}(F)
 =W^{-3-q}_{-\eps,b}(\Lambda^{2,q}\otimes F),
 \qquad q=0,1,2.
\end{equation}
Its cohomology is finite-dimensional. For $F=\OO_Y$ its degree-one
cohomology is zero. In particular every closed scalar $(2,1)$-current
in this space which is zero on a positive tail is exact.

If a degree-zero current
$\beta\in W^{-3}_{-\eps,b}(\Lambda^{2,0}\otimes F)$ satisfies
$\dbar_F\beta=J$ for a degree-one current $J$ supported in
$\{h\le C\}$, then $\beta$ vanishes on a positive tail. Finally, an
order-zero $(2,1)$-current supported in $\{h\le C\}$ whose masses on
negative unit height slabs are at most $C_0e^{\mu l}$ belongs to this
complex whenever $b>\mu$. The same Sobolev order and pair of rates can
be used for any family with that common exponent $\mu$, regardless of
its individual constants and upper support bounds.
\end{theorem}

\begin{proof}
Finite-dimensionality is Proposition~\ref{weighted:fredholm} with
$s=-3$. We first give the direct exactness argument needed for
bounded-above support. If $J$ is scalar and zero on the positive tail,
it belongs to the single-linear-rate complex $(-b,b)$: its negative
weight is unchanged and its positive part has compact support. That
complex is acyclic. Its primitive has rates $(-b,b)$ and therefore also
belongs to $(-\eps,b)$, since $b>\eps$.

The full scalar degree-one vanishing can also be checked without a
contour argument in higher degree. Keep $b$ fixed and continuously vary
the positive rate from $-b$ to $-\eps$. After conjugation the even-to-odd
Hodge operators have the same domain $W^1$ and range $L^2$, and vary
continuously in operator norm: only bounded order-zero coefficients
depending on the derivative of the weight change. They remain
Fredholm by Proposition~\ref{weighted:fredholm}, since the positive
rate stays negative. Their index is consequently constant, and is
zero at the acyclic single-linear pair $(-b,b)$. At every point of
this path, $H^{2,0}=0$ by positive-tail uniqueness. By weighted duality
$H^{2,2}$ is dual to the scalar holomorphic functions with a positive
rate at the positive end and the negative rate $-b$ at the negative
end. Every such global scalar function is constant: subtract its
positive-end coefficient and use Proposition~\ref{weighted:tail}.
The negative-end decay requirement excludes a nonzero constant.
Thus $H^{2,2}=0$ as well. Index zero now implies $H^{2,1}=0$.

For the primitive assertion, on $\{h>C\}$ the distributional section
$\beta$ solves the degree-zero equation for $K_Y\otimes F$. Local
elliptic regularity makes it smooth there. Uniform interior estimates
on shorter tails turn its weighted $W^{-3}$ bound into finite
exponential growth in every positive Sobolev order. Proposition
\ref{weighted:tail} therefore makes it zero on a positive tail. This
holds for every primitive in the stated space.

To verify the last membership assertion, localize to the uniformly
bounded coordinate charts on one enlarged slab. In real dimension
four, $W^4$ embeds continuously into $C^0$, with a uniform constant on
these charts. An order-zero current of mass $M$ thus has $W^{-4}$ norm
at most $CM$. A uniformly locally finite partition of unity and the
bounded derivatives of the weight give an upper bound for the squared
weighted norm by a constant multiple of
\[
 \sum_{l\ge0} e^{-2bl}(C_0e^{\mu l})^2
 =C_0^2\sum_{l\ge0}e^{-2(b-\mu)l}<\infty.
\]
There are only finitely many additional slabs between the zero level
and the upper support bound, and none beyond it. This proves
membership in $W^{-4}_{-\eps,b}$, which is degree one
of~\eqref{weighted:current-complex}. The constants of individual
currents enter their norms, not the choice of spaces or the dimension
of their common cohomology group.
\end{proof}

In particular, if an order-zero $(2,1)$-current $C$ with coefficients
in $F\in\{\OO,E\}$ is bounded above in height and satisfies
$\Mass(C;Y_l)\le C_Ce^{\mu l}$, then for any $\eps>0$ and
$b>\max\{\mu,\eps\}$ the theorem places it in
\[
 W^{-4}_{-\eps,b}(\Lambda^{2,1}\otimes F).
\]
The inequality $b>\mu$ makes the weighted mass sum converge. The base
Sobolev order remains $-3$ and the degree-one order remains $-4$,
independently of both the current and its mass exponent. We will always
fix the rates before comparing current classes; no dimension bound
uniform over changing weights is needed.

\subsection{Residue independence}

If a linear combination of our currents is exact, its primitive
vanishes first on a positive tail and then off the supporting curves
by analytic continuation. A local distribution supported on a curve
cannot have the required order-zero residue as its Dolbeault
derivative. This gives the independence statement that we will use
both for scalar and for bundle-valued differentials.

\begin{proposition}[Residue obstruction]\label{curve:residue}
Fix $F\in\{\OO,E\}$ and rates $(-\eps,b)$ as in
Theorem~\ref{weighted:currents}.  For distinct closed irreducible curves
$B_i\subset Y$ bounded above in height, let $\eta_i$ be nonzero
holomorphic differentials on their normalizations with coefficients
in $f_i^*F$.  Suppose their pushforward currents have negative-end mass
exponents strictly smaller than $b$.  Then the classes
\[
 [C_{\eta_i}]\in H^{2,1}_{-\eps,b}(Y,F)
\]
are linearly independent: every finite subfamily is independent.
For $F=\OO$, even one such differential is impossible.
\end{proposition}

\begin{proof}
The mass criterion in Theorem~\ref{weighted:currents} puts the currents
in the indicated complex.
Suppose a finite linear combination is exact:
\[
 \sum_{i=1}^r a_i C_{\eta_i}=\dbar_F\beta.
\]
The sum is zero on a positive tail.  The primitive assertion in
Theorem~\ref{weighted:currents} says that $\beta$ is zero there as well.
On $Y\setminus\bigcup_iB_i$, its distributional Dolbeault derivative
vanishes.  Local elliptic regularity makes it a holomorphic section of
$K_Y\otimes F$ on that complement.  The complement is connected: a
path in the connected real four-manifold can be perturbed off the
smooth real-codimension-two strata of a finite union of analytic
curves, and off its locally discrete singular strata.  Only finitely
many charts are needed along the compact path.  Analytic continuation
from the positive tail therefore gives
\[
 \supp\beta\subset\bigcup_{i=1}^rB_i.
\]

Suppose $a_i\ne0$.  Choose a smooth point of $B_i$, outside the other
curves, where $\eta_i\ne0$.  Such a point exists since the singular
and intersection sets are locally discrete and a nonzero holomorphic
differential has isolated zeros.  In holomorphic coordinates $(z,w)$
and a holomorphic frame of $F$, the curve is $w=0$, the differential is
$a(z)\,dz$ with $a\not\equiv0$, and
$\beta=u\,dz\wedge dw$ for a distribution supported on $w=0$.
Its local finite normal-jet expansion has the form
\begin{equation}\label{curve:normal-jets}
 u=\sum_{p+q\le N}u_{pq}(z)\,
       \partial_w^p\partial_{\bar w}^q\delta_0(w),
\end{equation}
with distributional tangential coefficients.  This is the support
theorem for distributions \cite[Theorem~2.3.5]{Hormander1990}; locally it also follows
by taking a finite-order bound for $u$, Taylor-expanding a test function
in its normal variables, and cutting off in $|w|<\delta$.  A test
function whose normal derivatives through order $N$ vanish on $w=0$
has $C^N$ norm $O(\delta)$ after that cutoff, and hence pairs to zero
as $\delta\to0$.  Thus only the displayed finite Taylor jet can matter.
The normal jets in \eqref{curve:normal-jets} are independent, as is
seen by prescribing those jets of test functions.

The coefficient of the normal exterior-form direction in
$\dbar\beta$ is $\partial_{\bar w}u$; its tangential direction cannot
contribute to that coefficient.  The corresponding coefficient of the
right-hand side is a nonzero constant multiple of
$a_i a(z)\delta_0(w)$.  But normal differentiation in
\eqref{curve:normal-jets} increases $q$ by one.  Comparing the largest
$q+1$, then descending, forces every $u_{pq}$ to vanish and leaves
zero in place of this nonzero order-zero delta mass.  This contradiction
proves independence.  Notice that the familiar principal-value
meromorphic primitive of a residue is not supported on the divisor;
the support conclusion above is essential.

For $F=\OO$, Theorem~\ref{weighted:currents} makes every closed current
of the kind under consideration exact: its support is bounded above,
so a single linear weight $(-b,b)$ supplies the scalar primitive.
The independence assertion then rules out even one nonzero current.
\end{proof}

\subsection{The topology of the normalization}

An $L^2$ differential satisfies the required current mass bound by
Cauchy--Schwarz and the area estimate. Thus the scalar residue
obstruction excludes such differentials. The next proposition turns
this analytic conclusion into a restriction on the topology of the
normalization, including when it is noncompact.

\begin{proposition}[Genus zero]\label{curve:genus}
The normalization $R$ of any component of $D$ has no nonzero
square-integrable scalar holomorphic differential and has genus zero.
If $R$ is compact, it is biholomorphic to $\PP^1$.
If it is noncompact, it is biholomorphic to a domain in $\PP^1$,
without a restriction on the connectivity of that domain.
\end{proposition}

\begin{proof}
Let $\eta$ be a square-integrable holomorphic differential on $R$.
Cauchy--Schwarz in \eqref{curve:mass-estimate} gives
\[
 \Mass(C_\eta;Y_l)
 \le c\|\eta\|_{L^2(R)}
       \left(\int_{f^{-1}(Y_l)}|df|^2dA_R\right)^{1/2}
 \le C\|\eta\|_{L^2(R)}e^{Al/2}.
\]
Here the energy of the holomorphic normalization is a fixed multiple
of the reduced curve area; ramification and singular points are
measure-zero sets.  Lemma~\ref{curve:area} gives the last bound.
The scalar case of Proposition~\ref{curve:residue} forces $\eta=0$.

We give the topological consequence also for open $R$.  A handle in an
oriented surface contains two embedded circles with algebraic
intersection one.  Smooth closed forms $\alpha,\beta$, supported in
small annular neighborhoods of these circles and dual to them, have
\[
 \int_R\alpha\wedge\beta=1.
\]
For example a bump function of the transverse coordinate, times its
differential and with integral one, gives each such dual form.
Choose any smooth conformal metric and let $V$ be the closure of
$dC_c^\infty(R,\RR)$ in the real Hilbert space of $L^2$ one-forms.
Set $u=\alpha-\operatorname{proj}_V\alpha$.  Forms in $V$ are weakly
closed, so $du=0$; orthogonality to $dC_c^\infty$ gives $d^*u=0$.
Integration by parts against the compactly supported closed form
$\beta$, followed by $L^2$ continuity, gives
\[
 \int_Ru\wedge\beta=\int_R\alpha\wedge\beta=1.
\]
Thus $u\ne0$.  Local elliptic regularity makes $u$ smooth.
With the complex orientation, $u+i*u$ is a nonzero $(1,0)$-form;
the two equations for $u$ make it closed and hence holomorphic.
It is $L^2$, a contradiction.  No closed-range assertion for the
de Rham differential on a noncompact surface has been used.

There is consequently no handle in any compact subsurface of $R$.
This is genus zero.  Equivalently, every Jordan curve separates:
a nonseparating Jordan curve can be joined from one side to the other
in its complement, producing a second circle with intersection one
and hence a handle in a compact neighborhood.  A compact genus-zero
Riemann surface is $\PP^1$.  For noncompact $R$, apply Koebe's general
planar uniformization theorem in the form
\cite[Theorem~1.1]{Hemasundar2011}: every planar Riemann surface admits
a conformal embedding into the Riemann sphere.  The criterion that every
Jordan curve separates is its planarity hypothesis.  The theorem allows
arbitrary connectivity, so no finite-type assumption on $R$
is being added.  Its image is open, and thus is the asserted domain.
\end{proof}
